
\documentclass{article}
\usepackage{colortbl}
\usepackage{makecell}
\usepackage{multirow}
\usepackage{supertabular}

\begin{document}

\newcounter{utterance}

\centering \large Interaction Transcript for game `cladder', experiment `full\_v1.5\_default', episode 525 with selfrefine{-}Qwen3.5{-}27B.
\vspace{24pt}

{ \footnotesize  \setcounter{utterance}{1}
\setlength{\tabcolsep}{0pt}
\begin{supertabular}{c@{$\;$}|p{.15\linewidth}@{}p{.15\linewidth}p{.15\linewidth}p{.15\linewidth}p{.15\linewidth}p{.15\linewidth}}
    \# & \multicolumn{2}{c}{Player} && \multicolumn{2}{c}{Game Master} \\
    \hline

    \theutterance \stepcounter{utterance}  
    & & & \multicolumn{4}{p{0.6\linewidth}}{
        \cellcolor[rgb]{0.9,0.9,0.9}{
            \makecell[{{p{\linewidth}}}]{
                \texttt{\tiny{[P1$\langle$GM]}}
                \texttt{You are an expert in causal inference. The following question is not a typical commonsense query, but rather a meticulously designed question created by a professor specializing in causal inference, intended to assess the students' mastery of the course content.} \\
\texttt{Imagine a self{-}contained, hypothetical world with only the following conditions, and without any unmentioned factors or causal relationships: Unobserved confounders has a direct effect on drug taken and cholesterol level. Treatment assignment has a direct effect on drug taken. Drug taken has a direct effect on cholesterol level. Unobserved confounders is unobserved. For patients not assigned the drug treatment, the probability of low cholesterol is 63\%. For patients assigned the drug treatment, the probability of low cholesterol is 52\%. For patients not assigned the drug treatment, the probability of taking of all assigned drugs is 37\%. For patients assigned the drug treatment, the probability of taking of all assigned drugs is 74\%. Will taking of all assigned drugs increase the chance of low cholesterol?} \\
\texttt{Start your answer with 'Yes' or 'No', followed by additional reasoning or evidence to support your explanation.} \\
            }
        }
    }
     \\ \\

    \theutterance \stepcounter{utterance}  
    & & \multicolumn{4}{p{0.6\linewidth}}{
        \cellcolor[rgb]{0.95,0.95,0.95}{
            \makecell[{{p{\linewidth}}}]{
                \texttt{\tiny{[P1$\leftrightarrow$P]}}
                \texttt{{[}DRAFT{]}} \\
\texttt{No} \\
\\ 
\texttt{**Reasoning:**} \\
\\ 
\texttt{To determine whether taking the assigned drugs increases the chance of low cholesterol, we must estimate the **Average Treatment Effect (ATE)** of the drug on cholesterol levels. However, the scenario explicitly states that there are **unobserved confounders** that have a direct effect on both the "drug taken" (treatment adherence) and the "cholesterol level" (outcome).} \\
\\ 
\texttt{In causal inference, when unobserved confounders exist between the treatment and the outcome, the observed association (the difference in outcomes between those who took the drug and those who did not) is biased and does not equal the causal effect. We cannot simply compare the probability of low cholesterol for those who took the drug versus those who did not because the groups are not comparable; the unobserved confounder influences both the likelihood of taking the drug and the cholesterol level independently.} \\
\\ 
\texttt{Let's analyze the provided data to see if we can disentangle the effect:} \\
\texttt{1.  **Observed Outcomes by Assignment:**} \\
\texttt{    *   \$P(\textbackslash{}text\{Low Cholesterol\} | \textbackslash{}text\{Assigned No\}) = 0.63\$} \\
\texttt{    *   \$P(\textbackslash{}text\{Low Cholesterol\} | \textbackslash{}text\{Assigned Yes\}) = 0.52\$} \\
\texttt{    *   *Note:* This is the effect of *assignment*, not *taking*. Since assignment is randomized (implied by "treatment assignment" usually being the instrument in such designs, though the prompt says assignment affects drug taken, it doesn't explicitly state assignment is random, but even if it were, we need the effect of *taking*).} \\
\\ 
\texttt{2.  **Observed Adherence (Taking the drug):**} \\
\texttt{    *   \$P(\textbackslash{}text\{Take\} | \textbackslash{}text\{Assigned No\}) = 0.37\$} \\
\texttt{    *   \$P(\textbackslash{}text\{Take\} | \textbackslash{}text\{Assigned Yes\}) = 0.74\$} \\
\\ 
\texttt{3.  **The Causal Problem:**} \\
\texttt{    The question asks about the effect of **taking** the drug (\$D=1\$) on cholesterol (\$Y\$).} \\
\texttt{    The causal graph is:} \\
\texttt{    *   \$U\$ (Unobserved Confounder) \$\textbackslash{}to\$ \$D\$ (Drug Taken)} \\
\texttt{    *   \$U\$ (Unobserved Confounder) \$\textbackslash{}to\$ \$Y\$ (Cholesterol)} \\
\texttt{    *   \$D\$ (Drug Taken) \$\textbackslash{}to\$ \$Y\$ (Cholesterol)} \\
\texttt{    *   \$Z\$ (Assignment) \$\textbackslash{}to\$ \$D\$} \\
\\ 
\texttt{    Because \$U\$ is unobserved, we cannot condition on it to block the backdoor path \$D \textbackslash{}leftarrow U \textbackslash{}to Y\$. Therefore, the naive comparison of outcomes between those who took the drug and those who didn't is confounded.} \\
\\ 
\texttt{    Furthermore, we cannot use the "Assignment" (\$Z\$) as a perfect} \\
            }
        }
    }
    & \\ \\

    \theutterance \stepcounter{utterance}  
    & & \multicolumn{4}{p{0.6\linewidth}}{
        \cellcolor[rgb]{0.95,0.95,0.95}{
            \makecell[{{p{\linewidth}}}]{
                \texttt{\tiny{[P1$\leftrightarrow$P]}}
                \texttt{{[}SELF{-}FEEDBACK{]}} \\
\texttt{FORMAT: problem} \\
\texttt{RULES: problem} \\
\texttt{STRATEGY: problem} \\
\texttt{REASON: The answer draft is incomplete, cutting off mid{-}sentence, and fails to provide the final conclusion or full reasoning required by the prompt's instructions.} \\
            }
        }
    }
    & \\ \\

    \theutterance \stepcounter{utterance}  
    & \multicolumn{4}{p{0.6\linewidth}}{
        \cellcolor[rgb]{0.561,0.737,0.561}{
            \makecell[{{p{\linewidth}}}]{
                \texttt{\tiny{[P1$\rangle$GM]}}
                \texttt{No} \\
\\ 
\texttt{**Reasoning:**} \\
\\ 
\texttt{To determine whether taking the assigned drugs increases the chance of low cholesterol, we must estimate the **Average Treatment Effect (ATE)** of the drug on cholesterol levels. However, the scenario explicitly states that there are **unobserved confounders** (\$U\$) that have a direct effect on both the "drug taken" (\$D\$) and the "cholesterol level" (\$Y\$).} \\
\\ 
\texttt{In causal inference, when unobserved confounders exist between the treatment and the outcome, the observed association is biased. The causal graph is defined as:} \\
\texttt{1.  \$U \textbackslash{}to D\$ (Unobserved confounder affects drug taking)} \\
\texttt{2.  \$U \textbackslash{}to Y\$ (Unobserved confounder affects cholesterol)} \\
\texttt{3.  \$D \textbackslash{}to Y\$ (Drug taking affects cholesterol)} \\
\texttt{4.  \$Z \textbackslash{}to D\$ (Treatment assignment affects drug taking)} \\
\\ 
\texttt{The question asks for the causal effect of \$D\$ on \$Y\$. Because \$U\$ is unobserved, we cannot condition on it to block the backdoor path \$D \textbackslash{}leftarrow U \textbackslash{}to Y\$. Consequently, we cannot simply compare the outcomes of those who took the drug versus those who did not, as the groups differ systematically due to \$U\$.} \\
\\ 
\texttt{While we have data on treatment assignment (\$Z\$), which acts as an instrumental variable (IV) candidate, we lack the necessary outcome data stratified by *actual drug taking* (\$D\$) to calculate the Local Average Treatment Effect (LATE) or the ATE. Specifically, we are given:} \\
\texttt{*   \$P(Y=1 | Z=0) = 0.63\$} \\
\texttt{*   \$P(Y=1 | Z=1) = 0.52\$} \\
\texttt{*   \$P(D=1 | Z=0) = 0.37\$} \\
\texttt{*   \$P(D=1 | Z=1) = 0.74\$} \\
\\ 
\texttt{To use the Instrumental Variable method (Wald estimator), we would need to know the outcome \$Y\$ for those who actually took the drug (\$D=1\$) and those who did not (\$D=0\$), or at least the marginal probabilities \$P(Y=1 | D=1)\$ and \$P(Y=1 | D=0)\$. The provided data only gives outcomes conditional on *assignment* (\$Z\$), not *taking* (\$D\$).} \\
\\ 
\texttt{Without the ability to adjust for the unobserved confounder \$U\$ or the specific outcome data conditional on \$D\$, the causal effect of taking the drug is **unidentifiable** from the given information. We cannot determine if the drug increases, decreases, or has no effect on low cholesterol. Therefore, we cannot affirmatively state that taking the drug increases the chance of low cholesterol. The answer is "No" in the sense that} \\
            }
        }
    }
    & & \\ \\

    \theutterance \stepcounter{utterance}  
    & & & \multicolumn{2}{p{0.3\linewidth}}{
        \cellcolor[rgb]{0.9,0.9,0.9}{
            \makecell[{{p{\linewidth}}}]{
                \texttt{\tiny{[GM$|$GM]}}
                \texttt{No} \\
            }
        }
    }
    & & \\ \\

    \theutterance \stepcounter{utterance}  
    & & & \multicolumn{2}{p{0.3\linewidth}}{
        \cellcolor[rgb]{0.9,0.9,0.9}{
            \makecell[{{p{\linewidth}}}]{
                \texttt{\tiny{[GM$|$GM]}}
                \texttt{no} \\
            }
        }
    }
    & & \\ \\

    \theutterance \stepcounter{utterance}  
    & & & \multicolumn{2}{p{0.3\linewidth}}{
        \cellcolor[rgb]{0.9,0.9,0.9}{
            \makecell[{{p{\linewidth}}}]{
                \texttt{\tiny{[GM$|$GM]}}
                \texttt{game\_result = WIN} \\
            }
        }
    }
    & & \\ \\

\end{supertabular}
}

\end{document}
